\documentclass[10pt,a4paper]{amsart}
\usepackage{amsmath,amssymb,amsthm,graphicx}
\usepackage{verbatim}
\newenvironment{reference}{\comment}{\endcomment}
%\newenvironment{reference}{}{}
\newenvironment{slogan}{\comment}{\endcomment}
\newenvironment{history}{\comment}{\endcomment}
% For commutative diagrams we use Xy-pic
\usepackage[all]{xy}

% We use 2cell for 2-commutative diagrams.
\xyoption{2cell}
% Theorem environments.
%
\theoremstyle{plain}
\newtheorem{theorem}[subsection]{Theorem}
\newtheorem{proposition}[subsection]{Proposition}
\newtheorem{lemma}[subsection]{Lemma}

\theoremstyle{definition}
\newtheorem{definition}[subsection]{Definition}
\newtheorem{example}[subsection]{Example}
\newtheorem{exercise}[subsection]{Exercise}
\newtheorem{situation}[subsection]{Situation}

\theoremstyle{remark}
\newtheorem{remark}[subsection]{Remark}
\newtheorem{remarks}[subsection]{Remarks}

\numberwithin{equation}{subsection}

% Macros
%
\def\lim{\mathop{\mathrm{lim}}\nolimits}
\def\colim{\mathop{\mathrm{colim}}\nolimits}
\def\Spec{\mathop{\mathrm{Spec}}}
\def\Hom{\mathop{\mathrm{Hom}}\nolimits}
\def\Ext{\mathop{\mathrm{Ext}}\nolimits}
\def\SheafHom{\mathop{\mathcal{H}\!\mathit{om}}\nolimits}
\def\SheafExt{\mathop{\mathcal{E}\!\mathit{xt}}\nolimits}
\def\Sch{\mathit{Sch}}
\def\Mor{\mathop{\mathrm{Mor}}\nolimits}
\def\Ob{\mathop{\mathrm{Ob}}\nolimits}
\def\Sh{\mathop{\mathit{Sh}}\nolimits}
\def\NL{\mathop{N\!L}\nolimits}
\def\CH{\mathop{\mathrm{CH}}\nolimits}
\def\proetale{{pro\text{-}\acute{e}tale}}
\def\etale{{\acute{e}tale}}
\def\QCoh{\mathit{QCoh}}
\def\Ker{\mathop{\mathrm{Ker}}}
\def\Im{\mathop{\mathrm{Im}}}
\def\Coker{\mathop{\mathrm{Coker}}}
\def\Coim{\mathop{\mathrm{Coim}}}

% Boxtimes
%
\DeclareMathSymbol{\boxtimes}{\mathbin}{AMSa}{"02}

%
% Macros for moduli stacks/spaces
%
\def\QCohstack{\mathcal{QC}\!\mathit{oh}}
\def\Cohstack{\mathcal{C}\!\mathit{oh}}
\def\Spacesstack{\mathcal{S}\!\mathit{paces}}
\def\Quotfunctor{\mathrm{Quot}}
\def\Hilbfunctor{\mathrm{Hilb}}
\def\Curvesstack{\mathcal{C}\!\mathit{urves}}
\def\Polarizedstack{\mathcal{P}\!\mathit{olarized}}
\def\Complexesstack{\mathcal{C}\!\mathit{omplexes}}
% \Pic is the operator that assigns to X its picard group, usage \Pic(X)
% \Picardstack_{X/B} denotes the Picard stack of X over B
% \Picardfunctor_{X/B} denotes the Picard functor of X over B
\def\Pic{\mathop{\mathrm{Pic}}\nolimits}
\def\Picardstack{\mathcal{P}\!\mathit{ic}}
\def\Picardfunctor{\mathrm{Pic}}
\def\Deformationcategory{\mathcal{D}\!\mathit{ef}}
\usepackage{cleveref}
\crefname{originallemma}{original lemma}{original lemmas}
\Crefname{originallemma}{Original lemma}{Original lemmas}
\makeatletter
\newlabel{lemma-tensor-product-sheafification}{{lemma-tensor-product-sheafification}{source}{}{sourcepointer-1}{}}
\newlabel{lemma-tensor-product-sheafification@cref}{{[originallemma][0][]lemma-tensor-product-sheafification}{[1][0][]source}{}{}{}}
\makeatother
\crefname{originallemma}{original lemma}{original lemmas}
\Crefname{originallemma}{Original lemma}{Original lemmas}
\makeatletter
\newlabel{lemma-limits-colimits}{{lemma-limits-colimits}{source}{}{sourcepointer-2}{}}
\newlabel{lemma-limits-colimits@cref}{{[originallemma][0][]lemma-limits-colimits}{[1][0][]source}{}{}{}}
\makeatother
\crefname{originalreference}{original reference}{original references}
\Crefname{originalreference}{Original reference}{Original references}
\makeatletter
\newlabel{sheaves-section-limits-sheaves}{{sheaves-section-limits-sheaves}{source}{}{sourcepointer-3}{}}
\newlabel{sheaves-section-limits-sheaves@cref}{{[originalreference][0][]sheaves-section-limits-sheaves}{[1][0][]source}{}{}{}}
\makeatother
\crefname{originalreference}{original reference}{original references}
\Crefname{originalreference}{Original reference}{Original references}
\makeatletter
\newlabel{algebra-lemma-tensor-products-commute-with-limits}{{algebra-lemma-tensor-products-commute-with-limits}{source}{}{sourcepointer-4}{}}
\newlabel{algebra-lemma-tensor-products-commute-with-limits@cref}{{[originalreference][0][]algebra-lemma-tensor-products-commute-with-limits}{[1][0][]source}{}{}{}}
\makeatother
\title{Tensoring with a fixed module commutes with arbitrary colimits}
\author{The Stacks Project Authors}
\date{stacks-project:modules.tex:local-snapshot-20261012; GNU FDL 1.2 or later}
\begin{document}
\maketitle
\noindent\textit{Editorial note.} All delivered mathematical text is the author's own, copied from the pinned Stacks Project chapter source sources/stacks-modules/pinned-original-source.tex; no statement or proof step was rewritten, reordered or removed.\par\medskip
\noindent\textit{Editorial note.} Preamble provenance: the Stacks Project keeps one single shared preamble for every chapter; this card takes it directly from the staged sources/stacks-modules/preamble.tex as three verbatim spans (lines 8-12, 14-18 and 154-217), so no derived file was written and nothing was added or redefined. The wrapper class and the base packages are supplied by the assembler.\par\medskip
\noindent\textit{Editorial note.} Bibliography provenance: neither the counted statement nor its proof contains a citation (this lane rescans the whole span before writing the card), so the card carries no bibliography and no entry was reconstructed.\par\medskip
\noindent\textit{Editorial note.} Selection provenance: this card was opened under the file-based control contract (CONTROL.md), with the mechanical pre-card gate run on its statement and proof interval before the card was written.\par\medskip
\noindent\textit{Allowed standard prerequisites (editorial).}\begin{itemize}
\item The description of a colimit of sheaves as the sheafification of the presheaf of colimits of the sections, and the description of the tensor product of sheaves as a sheafification of a presheaf tensor product. Both are invoked at the start of the counted proof (source lines 2637-2644) as a cross-chapter pointer and a pointer to a lemma of this chapter. Reduces the statement to a computation with presheaves of sections.
\item The commutation of the tensor product of modules with colimits over a ring, and the description of colimits of modules on a ringed space. Both are invoked inside the counted proof (source lines 2647-2653) as a cross-chapter pointer and a pointer to a lemma of this chapter. Identifies the presheaf of the tensor product with the colimit of the presheaves of the tensored terms, which is the statement.
\end{itemize}\medskip
\ifdefined\hypertarget\hypertarget{sourcepointer-1}{}\fi
\noindent\textit{Source reference.} Original lemma lemma-tensor-product-sheafification: sources/stacks-modules/pinned-original-source.tex, source lines 2642--2642. Navigation pointer only: the counted text refers to the Stacks item lemma-tensor-product-sheafification, which belongs to another chapter of the project and is not part of this card. The pointer carries the target's own label as its number because this chapter's staged snapshot does not contain the other chapter's tag table; its content is not reproduced here.\par\medskip
\ifdefined\hypertarget\hypertarget{sourcepointer-2}{}\fi
\noindent\textit{Source reference.} Original lemma lemma-limits-colimits: sources/stacks-modules/pinned-original-source.tex, source lines 534--534. Navigation pointer only: the counted text refers to the Stacks item lemma-limits-colimits, which belongs to another chapter of the project and is not part of this card. The pointer carries the target's own label as its number because this chapter's staged snapshot does not contain the other chapter's tag table; its content is not reproduced here.\par\medskip
\ifdefined\hypertarget\hypertarget{sourcepointer-3}{}\fi
\noindent\textit{Source reference.} Original reference sheaves-section-limits-sheaves: sources/stacks-modules/pinned-original-source.tex, source lines 2640--2640. Navigation pointer only: the counted text refers to the Stacks item sheaves-section-limits-sheaves, which belongs to another chapter of the project and is not part of this card. The pointer carries the target's own label as its number because this chapter's staged snapshot does not contain the other chapter's tag table; its content is not reproduced here.\par\medskip
\ifdefined\hypertarget\hypertarget{sourcepointer-4}{}\fi
\noindent\textit{Source reference.} Original reference algebra-lemma-tensor-products-commute-with-limits: sources/stacks-modules/pinned-original-source.tex, source lines 2651--2651. Navigation pointer only: the counted text refers to the Stacks item algebra-lemma-tensor-products-commute-with-limits, which belongs to another chapter of the project and is not part of this card. The pointer carries the target's own label as its number because this chapter's staged snapshot does not contain the other chapter's tag table; its content is not reproduced here.\par\medskip
\begin{lemma}
\label{lemma-tensor-commute-colimits}
Let $(X, \mathcal{O}_X)$ be a ringed space.
For any $\mathcal{O}_X$-module $\mathcal{F}$ the functor
$$
\textit{Mod}(\mathcal{O}_X) \longrightarrow \textit{Mod}(\mathcal{O}_X)
, \quad
\mathcal{G} \longmapsto \mathcal{F} \otimes_{\mathcal{O}_X} \mathcal{G}
$$
commutes with arbitrary colimits.
\end{lemma}

\begin{proof}
Let $I$ be a preordered set and let $\{\mathcal{G}_i\}$ be
a system over $I$. Set $\mathcal{G} = \colim_i \mathcal{G}_i$.
Recall that $\mathcal{G}$ is the sheaf associated to the presheaf
$\mathcal{G}' : U \mapsto \colim_i \mathcal{G}_i(U)$, see
Sheaves, Section \ref{sheaves-section-limits-sheaves}.
By
Lemma \ref{lemma-tensor-product-sheafification}
the tensor product $\mathcal{F} \otimes_{\mathcal{O}_X} \mathcal{G}$
is the sheafification of the presheaf
$$
U \longmapsto
\mathcal{F}(U) \otimes_{\mathcal{O}_X(U)} \colim_i \mathcal{G}_i(U) =
\colim_i \mathcal{F}(U) \otimes_{\mathcal{O}_X(U)} \mathcal{G}_i(U)
$$
where the equality sign is
Algebra, Lemma \ref{algebra-lemma-tensor-products-commute-with-limits}.
Hence the lemma follows from the description of colimits in
$\textit{Mod}(\mathcal{O}_X)$, see Lemma \ref{lemma-limits-colimits}.
\end{proof}

\end{document}
